\documentclass[11pt]{article}
\usepackage[a4paper, margin=2.5cm]{geometry}
\usepackage[utf8]{inputenc}
\usepackage{lmodern}
\usepackage[T1]{fontenc}
\usepackage{amsmath, amssymb}
\usepackage{physics}

\title{Projectile Motion with Air Resistance}
\author{Ivan Pestana}
\date{\today}

\begin{document}
\maketitle

\section{Physical Model}
\subsection{Motivation}
When we consider a projectile launched from ground level with an initial
velocity \(v_0\) at an angle \(\theta\) above the horizontal, the motion can be
described by Newton's second law. If we ask the question of which angle
provides the maximum range, in the absence of air resistance, the answer is
famously \(45^\circ\). However this is not the case if we include air
resistance.

\subsection{Newton's Second Law}
Newton's second law for the projectile reads
\begin{equation}
\vb{F}_r + \vb{F}_g = m\vb{a},
\end{equation}
where $\vb{F}_g = (0, -mg)$ is the force of gravity. Writing this
equation component by component,
\begin{equation}
    a_x = \frac{F_{r,x}}{m}, \qquad
    a_y = -g - \frac{F_{r,y}}{m},
\end{equation}
where $F_{r,x}$ and $F_{r,y}$ are still unknown, since we have not yet
written $\vb{F}_r$ explicitly.

\subsection{The Drag Force}
When air resistance is included, the projectile is subject to a drag force
$\vb{F}_r$ that always acts in the opposite direction of the velocity $\vb{v}$.
Throughout this document, bold symbols denote vectors (e.g. $\vb{v}$),
while the same symbol in italics denotes its magnitude (e.g.
$v = |\vb{v}|$).

The form of the drag force depends on the Reynolds number,
\begin{equation}
\mathrm{Re} = \frac{Dv\rho}{\eta},
\end{equation}
where $D$ is the projectile's size, $\rho$ the density of air, and
$\eta$ its viscosity.

$\mathrm{Re}$ is the ratio of inertial to viscous forces in the
surrounding flow. If we have a low $\mathrm{Re}$, viscous forces
dominate and the drag force is linear,
\begin{equation}
    F_r \propto v \qquad \text{(linear, or Stokes, drag)},
\end{equation}
which is a regime typical of small objects moving slowly 
through a viscous fluid, such as a droplet. This regime holds for
$\mathrm{Re} \lesssim 1$; for example, a $10\ \mu\mathrm{m}$ fog
droplet settling in air at roughly $2.7\ \mathrm{cm/s}$ has
$\mathrm{Re} \approx 0.016$ \cite{lubarda2022}. Turning the
proportionality into an equality introduces a constant $b$, the linear
drag coefficient, which for a sphere is given by Stokes' law as
$b = 3\pi\eta D$. Since the force points opposite to the velocity,
\begin{equation}
\vb{F}_r = -b\,\vb{v}.
\end{equation}

On the other hand, if we have a high $\mathrm{Re}$, inertial forces
dominate and the drag force is quadratic,
\begin{equation}
    F_r \propto v^2 \qquad \text{(quadratic, or Newtonian, drag)},
\end{equation}
which is a regime typical of bigger objects moving through air,
such as a thrown ball or a fired projectile. This regime holds for
speeds between approximately $0.25\ \mathrm{m/s}$ and $53\ \mathrm{m/s}$,
corresponding to $10^3 < \mathrm{Re} < 2\times10^5$
\cite{chudinov2013}.

Between these two regimes, roughly $1 \lesssim \mathrm{Re} \lesssim 10^3$,
there is a middle zone where neither law works well, and the drag force
has no simple formula, instead being found from experiments or lookup
tables \cite{lubarda2022}. This middle zone does not matter here, since
projectile motion falls well inside the quadratic regime.

So far we only know that the drag force grows with the square of the
speed -- we do not yet have an exact formula. To turn this
proportionality into a real equation, we introduce a constant, $C$,
called the drag coefficient:
\begin{equation}
    F_r = C v^2.
\end{equation}
This constant packages together everything else that affects the
force: the projectile's shape, its cross-sectional area, and the
density of the air. In reality, $C$ can itself change slightly with
the speed of the projectile, $C = C(\vb{v})$. However, this variation
is small within the range of speeds we consider, so throughout this
document we treat $C$ as a fixed, known constant. This is an
approximation, but a reasonable one: it greatly simplifies the
equations while still capturing the essential physics of drag.

We still need the force as a vector, not just a size. To point it in
the right direction, we multiply by a unit vector -- a vector of
length one -- that points opposite to the motion. Dividing $\vb{v}$ by
its own length, $\vb{v}/v$, gives exactly that: it keeps the direction
of the velocity but has length one. So the drag force is
\begin{equation}
\vb{F}_r = -C\,v^2\,\frac{\vb{v}}{v} = -C\,v\,\vb{v},
\qquad v = |\vb{v}| = \sqrt{v_x^2+v_y^2},
\end{equation}
where the minus sign flips the direction, and one factor of $v$
cancels against the denominator.

Substituting this into the component equations above,
\begin{equation}
    a_x = -\frac{C}{m}\,v\,v_x, \qquad
    a_y = -g - \frac{C}{m}\,v\,v_y.
\end{equation}

\section{System of Differential Equations}

\section{Numerical Integration}

\section{Stopping Condition}

\section{How this turns into Code}

\begin{thebibliography}{9}

\bibitem{lubarda2022}
M.~V.~Lubarda, V.~A.~Lubarda,
\emph{A review of the analysis of wind-influenced projectile motion in the presence of linear and nonlinear drag force},
Archive of Applied Mechanics \textbf{92} (2022) 1997--2017.

\bibitem{chudinov2013}
P.~S.~Chudinov, V.~A.~Eltyshev, Y.~A.~Barykin,
\emph{Simple analytical description of projectile motion in a medium with quadratic drag force},
Latin American Journal of Physics Education \textbf{7} (2013).

\end{thebibliography}

\end{document}